\scanpage{29}

$k'$ est ``connexe'' i.e. pour tte telle décomposition de $1\in k'$, \inserted{(où $I\ne\varnothing$ et} tous les $e'_i$ sauf un sont nuls, l'autre $1$, \deleted{\unsure{donc}} (p. ex. $k'$ \deleted{\unsure{corps}} intègre, ou $k'/\mathrm{Nil}_{k'}$ intègre, \deleted{\unsure{ou plus}} ou \inserted{$k'$} un anneau local \ldots\ldots) alors les applications considérées sont bijectives.

% slanted
\margin{$*$ \uline{équivalent à la notion de connexité dans une catégorie} \ldots Exemple : $k'$ intègre}

Si $I=\varnothing$, on trouve un objet initial dans $\mathit{Aff}_k$

\[
\varnothing_k=V_0\ \text{(algèbre nulle sur $k$)}
\]

On a \deleted{\unsure{donc}}

\[
\begin{aligned}
\varnothing_k(k')&=\begin{cases}
\varnothing&\text{si }k'\ne0 \quad\text{i.e. }V_{k'}\nleftarrow\varnothing_k\\
\{e\}&\text{si }k'=0
\end{cases}\\
\rotatebox{90}{$=$}&\\
\mathrm{Hom}(V_{k'},\varnothing_k)&
\end{aligned}
\]

Donc c'est un \uline{objet initial strict}. On l'appelle le \uline{$k$-espace algébrique vide} (un \deleted{\unsure{léger|objet}} abus de langage \ldots,

$(*)$ Conoyaux $\mathrm{Coker}(u,v:R\rightrightarrows X)$ correspond à $\mathrm{Ker}(u,v:A\rightrightarrows B)$ mais n'a pas beaucoup de sens. O.K. si $R$ est une ``relation d'équivalence'' dans $X$ et si $R$ fini sur $X$ par $u$, i.e. $B$ une $A$-algèbre finie. On reviendra dessus avec la question du passage aux quotients.

Sous-espaces algébriques fermés d'un espace algébrique affine $\mathscr X\simeq V_A$ : ils correspondent aux \uline{idéaux} de $A$ avec relation d'ordre renversée

\textcircled{\unsure{11|10}} \uline{Spectre d'une $k$-algèbre $A$} : espace sous-jacent.

\[
A\longleftrightarrow\mathscr X=V_A
\]

On va associer un objet géométrique

\textcircled{a} \unsure{Elément|Ensemble} sous-jacent : ensemble \inserted{$X$} des ``lieux'' \inserted{(loci)} de ``pts géométriques'' de $\mathscr X$

\[
\Longleftrightarrow\quad\text{ensemble }\deleted{\unsure{de tous}}\inserted{\mathrm{Spec}(A)}\ \text{des idéaux premiers de } A
\]

% archive sheet number: 28
